\scanpage{6}
% archive foliation: 5 (upper right and lower left)
Alors
\[
G'=P\times^G G=\operatorname{ad}(P)
\]
En fait, il y a des cas intéressants où \deleted{\unsure{ne|on}} on peut trouver aucun foncteur fibre. Mais si $\mathcal C$ est de type fini, alors on \inserted{doit pouvoir} montrer qu'on peut trouver un foncteur fibre sur une extension finie de $K$.
\deleted{On \unsure{réussit} ainsi à décrire les}
\inserted{Quand il en est ainsi, \deleted{on doit \unsure{arriver}},}
admettant le résultat ``d'unicité'' de foncteurs fibres, \deleted{on doit} arrive à décrire $(\mathcal C,F)$ intrinsèquement en termes d'une gerbe sur un site convenable\dots

8. \uline{\unsure{Cas des catégories|Cas de catégorie} des motifs.}
\deleted{\uuline{lien avec le groupe fondamental profini.}}

$S$ schéma régulier. \inserted{connexe}
On définit $\mathcal{M}^+(S)$, $\mathcal{M}(S)$ en paraphrasant ce qui est fait sur un corps, en travaillant \unsure{les|avec} $X$ projectif lisse sur $S$, avec $C(X)\overset{\mathrm{def}}=C(X_\eta)$ ($\eta$ pt générique de $S$).
Ainsi, $\mathcal{M}(S)\unsure{-|\longrightarrow}\mathcal{M}(\eta)$ est pleinement fidèle, permettant d'interpréter $\mathcal{M}(S)$ comme ss-catégorie pleine de $\mathcal{M}(\eta)$.
De plus, pour $\forall$ morphisme $f:S'\longrightarrow S$ de schémas réguliers, on définit $f^*:\mathcal{M}(S)\longrightarrow \mathcal{M}(S')$ un $\otimes$-foncteur correspondant (c'est ici que sert l'hypothèse de \uline{régularité} sur les schémas\dots).
\margin{à expliciter\dots}

Tout point \inserted{géom.} \uline{complexe} de $S$ définit un foncteur fibre ``de Betti'' \deleted{$\unsure{F_{\mathrm{Betti}}=F_\xi}$}
\[
F_{\unsure{\xi|\zeta}}^{\mathrm{Betti}}=F_{\unsure{\xi|\zeta}}^B:
\mathcal{M}(S)\unsure{-|\longrightarrow}\operatorname{Modf}(\mathbb{Q}),
\]
d'où un groupe proalgébrique $G$ sur $\mathbb{Q}$, pro-réductif
